% =============================================================================
%  Sección 5: Despersonalización o Estandarización
% =============================================================================
\section{¿Despersonalización o Estandarización?}
\label{sec:despersonalizacion-estandarizacion}

En las secciones anteriores exploramos un paseo por las diferentes fases en las que la sociedad mexicana
ha tenido que experimentar en sus diferentes estratos; desde el gobierno cambiante de ideología y los grupos
generacionales de cada década. En este capítulo, exploraremos su relación social desde la perspectiva actual,
es decir; como es que la pérdida de diferentes aspectos que hacían única y diferenciable cada generación o grupo
de personas aportaba una diversidad y riqueza cultural. Riqueza que se ha ido perdiendo con el paso del tiempo.
Para ello se usarán diferentes fuentes que citan este problema a un nivel \termino{Macro - Teórico} con un
enfoque cualitativo.

\subsection{Definición de Conceptos}
\label{ssec:definicion-conceptos}

Para comprender las dinámicas contemporáneas de interacción mediada por plataformas y sistemas técnicos,
resulta imperativo establecer una distinción conceptual rigurosa entre dos fenómenos convergentes
pero ontológicamente distintos: la \termino{Despersonalización} y la \termino{Estandarización}.

\begin{itemize}
    \item \textbf{Despersonalización (Sociológica y Tecnológica):} Se define como el
            proceso mediante el cual el sujeto pierde su agencia e identidad cualitativa para ser
            reconfigurado como un perfil cuantificable o un conjunto de metadatos (\textit{dataificación}).
            Desde la sociología de la tecnología, implica la disolución de la empatía y la sustitución
            del vínculo humano por interfaces algorítmicas de gestión. El individuo deja de
            ser un ciudadano u ``otro'' con rostro para convertirse en un nodo anónimo,
            predecible y administrable dentro de las lógicas del ecosistema digital, diluyendo
            las particularidades humanas en favor del rendimiento sistémico \cite{han2014psicopolitica}.

    \item \textbf{Estandarización (Tecnológica):} En términos estrictamente técnicos e ingenieriles,
            se define como el proceso metódico y deliberado de establecer normas, protocolos y
            especificaciones uniformes para la construcción de hardware, software y redes.
            Su propósito estructural no interviene con la identidad humana, sino que busca garantizar
            la interoperabilidad, la escalabilidad y la reducción de la complejidad técnica.
            Es el mecanismo fundacional (como los protocolos de red o las normativas ISO/IEEE)
            que permite que arquitecturas digitales e infraestructuras heterogéneas puedan ensamblarse,
            comunicarse y operar sin fricción a escala global \cite{russell2014open}.

    \item \textbf{La Sociedad como Totalidad y el Hecho Social Digital (Émile Durkheim):} Metodológicamente,
            se concibe a la sociedad no como la mera suma de individuos, sino como una totalidad
            supraindividual (un ``todo'') que ejerce coacción sobre sus miembros a través de
            \textit{hechos sociales}.

    En el contexto de las plataformas, la estandarización tecnológica y la despersonalización operan como
    nuevos hechos sociales. El ecosistema digital impone sus protocolos y formas de interacción
    (métricas, algoritmos, perfiles) de manera coercitiva y exterior al individuo. La participación en
    estas plataformas se vuelve obligatoria para la integración social contemporánea, subordinando la
    acción individual a la estructura coercitiva de la totalidad de la red \cite{durkheim_reglas}.

    \item \textbf{Relaciones de Poder y el Panoptismo (Michel Foucault):} Se define al poder no como una
            posesión estática que se ejerce de arriba hacia abajo, sino como una red de relaciones
            inmanentes, microfísicas y omnipresentes que atraviesan todo el cuerpo social, produciendo
            saberes y disciplinando sujetos.

    Bajo este lente, la estandarización técnica de las plataformas es el andamiaje perfecto para
    la vigilancia panóptica moderna. La despersonalización de los usuarios en bases de datos permite
    que los algoritmos ejerzan una relación de poder invisible y asimétrica. Al categorizar, predecir
    y normalizar conductas, los sistemas técnicos no solo administran datos, sino que disciplinan y
    moldean la subjetividad del usuario, determinando qué es visible y qué es marginal
    \cite{foucault2002vigilar}.

    \item \textbf{Los Dispositivos y la Sociedad de Control (Gilles Deleuze):} Un dispositivo es una
            red heterogénea de discursos, instituciones, leyes y enunciados científicos que
            responden a una urgencia histórica. Deleuze advierte la transición de los encierros
            disciplinarios hacia \textit{sociedades de control}, operadas al aire libre mediante
            mecanismos de modulación continua.

    Las plataformas digitales contemporáneas son la materialización absoluta de estos dispositivos
    de control. A diferencia de las instituciones físicas del siglo pasado, el control digital es
    ondulatorio, operando mediante contraseñas, rastreo de ubicación y flujos de datos.
    La despersonalización algorítmica permite que el sujeto ya no sea confinado, sino
    continuamente rastreado y modulado en tiempo real, convirtiendo la estandarización del
    software en el mecanismo de control más ubicuo y expansivo \cite{deleuze1990postscriptum}.

    \item \textbf{Modernidad y Generación Líquida (Zygmunt Bauman):} Conceptualiza la fase actual
            de la modernidad como una etapa ``líquida'', caracterizada por la disolución de las
            estructuras sólidas (instituciones, empleos, identidades) en favor de la fluidez, la
            desregulación, la inestabilidad y el consumo efímero.

    Paradójicamente, la rigurosa estandarización técnica del internet es lo que hace posible la
    \textit{modernidad líquida}. Los usuarios se integran a una ``generación líquida''
    donde las interacciones interpersonales han sido despersonalizadas y sustituidas por conexiones de
    red: vínculos frágiles, desechables y fáciles de desconectar con un solo clic.
    La plataforma estandarizada ofrece la ilusión de pertenencia, mientras fomenta un entorno social
    donde el compromiso a largo plazo se disuelve en el consumo inmediato de perfiles e
    identidades \cite{bauman2003modernidad}.

\end{itemize}

\subsection{¿Identidad?}
\label{ssec:identidad}

El nombre oficial de México como nación es Estados Unidos Mexicanos, el cual está compuesto por 32 entidades federales llamados
estados; 31 de estos con el término \termino{Estados Libres y soberanos} y una \termino{Capital Federal} o
Ciudad de México. Aunque esto es legalmente cierto, en la práctica podemos separar al país por tres zonas
geográficas bien diferenciadas. La configuración territorial de la República Mexicana no responde únicamente a
límites político-administrativos, sino a una profunda fragmentación histórica, demográfica y económica que ha
moldeado identidades regionales fuertemente contrastantes. Partiendo de la regionalización geoeconómica de
\cite{bassols1992mexico} y el análisis de identidades de \cite{gimenez2007estudios}, es posible
estructurar el país en tres grandes macrorregiones, cada una con características socioculturales
y tensiones históricas particulares:

\begin{itemize}
    \item \textbf{Norte:} Comprende las entidades fronterizas y áridas
    (Baja California, Sonora, Chihuahua, Coahuila, Nuevo León y Tamaulipas) junto con
    estados como Sinaloa y Durango.
    \begin{itemize}
        \item \textit{Cultura y gastronomía:} Adaptada a las condiciones del desierto y la ganadería,
                predominan las tortillas de harina, la carne bovina asada, los alimentos
                deshidratados (machaca) y los lácteos.
        \item \textit{Acento:} Se caracteriza por una entonación recia, el acortamiento de
                palabras y una fuerte asimilación de anglicismos (\textit{espanglish}) dada la
                fricción constante con la frontera estadounidense.
        \item \textit{Religión y costumbres:} Aunque mayoritariamente católica, presenta un sincretismo
                indígena mucho menor que el resto del país. Impera una ética de trabajo ligada al
                individualismo, la colonización de territorios hostiles y la agroindustria.
        \item \textit{Tensiones (últimos 100 años):} Ha sido el epicentro del surgimiento y
                consolidación de la \textit{narcocultura} y la violencia sistémica transnacional.
                Asimismo, ha enfrentado crisis demográficas por la migración fronteriza,
                explotación estructural en la industria maquiladora y una histórica tensión
                por el estrés hídrico.
    \end{itemize}

    \item \textbf{Centro:} Abarca el Altiplano Central, el Bajío y parte del
    Occidente (Ciudad de México, Estado de México, Jalisco, Michoacán, Puebla,
    Querétaro, Guanajuato, Tlaxcala e Hidalgo).
    \begin{itemize}
        \item \textit{Cultura y gastronomía:} Representa la consolidación de la matriz
                mesoamericana y conventual colonial. Es la cuna de los platillos más
                elaborados (moles, tamales, pozole) fundamentados en la tríada del maíz, frijol y chile.
        \item \textit{Acento:} Varía desde el ``chilango'' capitalino hasta el ``tapatío'',
                caracterizados por cantinelas específicas y el uso arraigado de formulismos de
                cortesía y subordinación lingüística (el uso del ``mande'').
        \item \textit{Religión y costumbres:} Altamente conservador y católico, siendo el
                bastión del Guadalupanismo y del marianismo, con estructuras familiares
                tradicionales muy arraigadas.
        \item \textit{Tensiones (últimos 100 años):} Su problema central ha sido el
                macrocefalismo y la hiperurbanización. La extrema centralización del poder
                político y económico generó una aguda segregación socioespacial, colapso
                ecológico en cuencas (como el Valle de México) y fue el escenario de represión
                a los grandes movimientos sociales, obreros y estudiantiles de la segunda mitad
                del siglo XX.
    \end{itemize}

    \item \textbf{Sur:} Incluye los estados del Pacífico sur, la península y el sureste
    (Oaxaca, Chiapas, Guerrero, Veracruz, Tabasco, Campeche, Yucatán y Quintana Roo).
    \begin{itemize}
        \item \textit{Cultura y gastronomía:} Extraordinariamente diversa, fundamentada en
                ingredientes endémicos de la selva y el trópico. Destacan las técnicas
                precolombinas (enterrado, tlayudas, cochinita pibil) y la comida de mar en las costas.
        \item \textit{Acento:} Profundamente permeado por la fonética de las lenguas
                originarias (maya, zapoteco, mixteco) y, en las costas, por rasgos
                caribeños (aspiración de la letra ``s''). En zonas de Chiapas pervive el voseo.
        \item \textit{Religión y costumbres:} Existe un intenso sincretismo donde los
                ritos católicos coexisten orgánicamente con la cosmovisión y el
                calendario agrícola indígena. La organización social comunitaria
                (\textit{usos y costumbres}) es fundamental.
        \item \textit{Tensiones (últimos 100 años):} Ha sufrido un histórico abandono estatal
                y rezago en infraestructura. Su riqueza natural contrasta con la pobreza
                extrema de su población, derivando en conflictos por extractivismo,
                caciquismo y levantamientos armados (como el EZLN en 1994),
                evidenciando la fractura permanente entre el proyecto modernizador del
                Estado y lo que \cite{bonfil1987mexico} denominó el ``México profundo''.
    \end{itemize}
\end{itemize}

Todas estas diferencias, aunque claras y marcadas entre cada sección del país, comparten ciertas necesidades
dentro de un área que el urbanismo denomina \termino{Zonas metropolitanas}, áreas compuestas por dos o más
entidades que sobrepasan un tipo de gobierno, en su mayoría pueden ser estatal y cuando crecen más, pueden llegar
a ser federal. El caso más claro lo tenemos en la zona central con la conocida \zmvm. Aunque no es la única, se tienen
la Zona Metropolitana de Guadalajara en el estado de Jalisco y colindantes y la Zona Metropolitana de Monterrey en el
estado de Nuevo León. Las cuales han experimentado un creciente aumento en su población y economía en los últimos
años debido al aumento de importaciones, exportaciones, demanda de trabajo y demás bienes y servicios
\cite{inegi_enoe_laboral}.

Si es así \textbf{¿Por qué se dice que vivimos un proceso de estandarización/despersonalización en la personalidad?}:
La respuesta es más difícil de responder con datos cuantitativos, ya que la misma definición de personalidad a nivel
psicológico puede definirse como las conductas o acciones que tenemos con las personas que nos rodean y al constructo
que definen y moldean el comportamiento con nuestros semejantes dentro de un grupo conocido como sociedad. Esto deja mal
parado a este análisis, ya que sería difícil cuantificar un espectro con múltiples sombras y luces y aún así, podemos
irnos con la vertiente de otros textos. Según \cite{wirth_urbanismo} en la sección \termino{El urbanismo como modo de vida},
define que la Ciudad es un ente con su propio sistema que busca protegerse y sobrevivir a las necesidades del tiempo,
en su mismo capítulo, describe que una ciudad puede moldear el comportamiento de sus habitantes bajo las necesidades o
utilidades que requiera dicho sistema. El comportamiento de sus habitantes será un factor determinante en la supervivencia
de la misma. Lejos de pensar que esto es una decisión aislada o tomada por un grupo selecto de personas sentadas detrás de
un enorme escritorio en juntas directivas ultra secretas, responde a un sistema que se adapta al tiempo: La ciudad del siglo
XXI no opera bajo los mismos principios de las ciudades del siglo XX e incluso, bajo este enfoque cualitativo, podríamos decir
que dichos requerimientos no obedecen a ningún plan establecido para moldear el comportamiento de las masas a una planeación
estratégica. Pero si podemos notar ciertas fisuras y convergencias entre la sociedad actual; en los últimos años, los grupos
sociales emergentes de la \termino{Generación Z} parecen seguir dos tendencias claras: Una constante validación de su identidad
a partir de creencias, modas y bienes de consumo provenientes de las redes sociales más influyentes del momento. Esto no es algo
nuevo ni tampoco un \termino{Descubrimiento}, el ser humano siempre tiende a seguir un grupo social dominante por el mero
instinto de supervivencia social inherente a nuestra evolución. Sin embargo, esta réplica de ideas ya no obedece a figuras
que tengan una idea revolucionaria, disruptiva o intelectualmente poderosa. Más bien a quién tenga el poder económico en sus
manos, el poder mediático o simplemente quién goce de un momento de fama en el momento correcto. Este fenómeno ha sido
ampliamente documentado en la teoría social crítica. De acuerdo con Bauman, la dinámica del consumismo moderno altera
nuestra percepción del bienestar, lo que se retrata en la siguiente frase:
\citaclave{
  ``En el Mundo Moderno, la felicidad se encuentra al alcance
  de una tienda.''
  \citep{bauman2014mal}
}
En otras obras del mismo autor, se definen algunos conceptos clave como lo son \termino{La modernidad Líquida} o dicho de otra
manera, lo efímero del momento producido por las grandes multinacionales que controlan el poder económico y alteran de manera
significativa la sensación de bienestar de la población (en esta caso mexicana) a un espectro de bienestar que se
desvanece en cuanto lo tocas. Pero, aunque Bauman defina de manera magistral la \termino{Modernidad Líquida}, en un
país tan diverso como lo es México, el poder económico no es el único que apremia y reducir la culpa de otros sectores,
sería un error cualitativo. El país siempre se ha pintado como \termino{De raíces originarias} o con \termino{Cultura Ancestral},
aunque los mismos gobernantes no sepan que significa esto a cara del mundo. Muchas de las tradiciones \termino{milenarias},
son una amalgama y combinación de espectros políticos, ideológicos y culturalmente impuestos por actores que en su momento,
poco o nada tuvieron que ver con la formación de lo que hoy conocemos como \termino{región}. En su obra ensayística fundamental,
\textit{El laberinto de la soledad}, Octavio Paz elabora una aguda crítica sobre la psique colectiva y la
identidad del mexicano, definiéndola a través del hermetismo y el uso constante de ``máscaras'' sociales.
Para Paz, la sociedad mexicana se construye sobre un profundo sentimiento de orfandad histórica ---producto de la violenta
imposición de la Conquista y las contradicciones no resueltas de la posrevolución--- que obliga al individuo a cerrarse
ante el mundo exterior por temor a ser vulnerado o dominado (el arraigado miedo a ``rajarse''). Esta actitud defensiva y
de simulación constante, si bien protege la intimidad del sujeto ante un entorno que percibe como hostil, lo condena a
una soledad ontológica. En consecuencia, esta coraza cultural le impide establecer vínculos de comunión genuina con
el ``otro'', limitando su capacidad para reconciliarse con su pasado y enfrentar plenamente su modernidad
\cite{paz1959laberinto}. Para ejemplificar esto, usaremos uno de los casos más escondidos y menos explorados pero que
fácilmente puede ser replicado en cualquier \termino{Pueblo originario}, esto según la investigación de
\cite{lopezcaballero2017indigenas}. En el sur de la \zmvm existe un nivel de abandono latente por la diferencia geográfica
que dificulta el acceso de las autoridades, ocasionando que el \termino{estado} o \termino{gobierno}, solo sea llamado cuando
existe un problema que no se ha podido resolver entre los mismos pobladores, ocasionando un problema de legitimidad entre los
integrantes de la comunidad. En la obra \textit{Indígenas de la nación} se describe el origen ancestral de la alcaldía de Milpa
Alta, el cual no constituyó nunca un pueblo unido desde sus raíces prehispánicas, sino a una alianza surgida antes de la
conquista española en el año 1521; teniendo a 9 pueblos originarios que formaron los poblados actuales de la alcaldía:
\begin{table}[H]
    \centering
    \caption[Significado de los pueblos originarios de Milpa Alta]{Significado de los pueblos originarios de Milpa Alta.}
    \label{tab:pueblos_milpa_alta}
    \fontsize{9pt}{11pt}\selectfont % Ajuste de fuente para acomodar el texto denso
    \begin{tabularx}{\textwidth}{l l X X X}
        \toprule
        \textbf{Nombre actual} & \textbf{Forma en náhuatl} & \textbf{Análisis morfológico} & \textbf{Traducción literal} & \textbf{Significado interpretativo} \\
        \midrule
        Oztotepec & Oztōtēpēc & oztōtl (cueva) + tēpēc (en el cerro, sobre el cerro) & ``En el cerro de las cuevas'' & Lugar montañoso con grutas o cuevas naturales, posiblemente con función ritual o de refugio. \\
        \addlinespace
        Tlacoyucan & Tlācoyōcān & tlācoyōtl (figura alargada) + -cān (lugar de) & ``Lugar de tlacoyos'' o ``Lugar alargado'' & Puede aludir a la forma del terreno o a la producción del alimento tradicional. \\
        \addlinespace
        Tlacotenco & Tlacōtenco & tlacōtl (vara, bastón, caña) + -tenco (a la orilla de) & ``A la orilla de los carrizos'' & Sitio donde crecían cañas o varas junto a un borde o ribera. \\
        \addlinespace
        Tepenahuac & Tēpēnahuac & tēpētl (cerro) + nahuac (junto a, cerca de) & ``Junto al cerro'' & Asentamiento al pie o al costado de un cerro. \\
        \addlinespace
        Tecoxpa / Tecozpa & Tēcōxpan & tēcōxtli (piedra amarilla o miel) + -pan (sobre, en) & ``En la piedra amarilla'' & Lugar con piedra volcánica amarillenta o cantera. \\
        \addlinespace
        Micatlán / Miacatlán & Mīacatlān & mīa (mucho, abundante) + -tlān (lugar de) & ``Lugar abundante'' & Denota fertilidad o abundancia de recursos naturales. \\
        \addlinespace
        Atocpan & Ātōcpan & ātl (agua) + tocpān (sobre el lodo o barro) & ``Sobre el lodo'' o ``Lugar en el agua y el barro'' & Alude a suelos húmedos y fértiles, próximos a zonas lacustres. \\
        \addlinespace
        Ohtenco & Ōhtenco & ōhtli (camino) + -tenco (a la orilla de) & ``A la orilla del camino'' & Lugar ubicado junto a un sendero o vía principal; vinculado a rutas comerciales. \\
        \addlinespace
        Xicomulco & Xīcomōlco & xīcomōlli (jícaro, vasija) + -co (en) & ``En el jícaro'' o ``Lugar de jícaras'' & Sitio donde abundaban los árboles de jícaro o donde se fabricaban vasijas. \\
        \addlinespace
        Tecomitl & Tēcomitl & tēcomitl (cántaro, olla, vasija) & ``El cántaro'' o ``Lugar de cántaros'' & Hace referencia a producción o forma del terreno semejante a una vasija. \\
        \addlinespace
        Malacachtepec & Malacachtepēc & malacach-tli (trenza, cuerda enrollada) + tēpēc (en el cerro) & ``En el cerro de las trenzas'' o ``Lugar del cerro torcido'' & Puede referirse a colinas de forma ondulante o caminos sinuosos. \\
        \addlinespace
        Cuauhtenco & Cuāuhtenco & cuāuhtli (águila o árbol) + -tenco (a la orilla de) & ``A la orilla del árbol / del bosque / del águila'' & Sitio ubicado al borde de un bosque o asociado al símbolo del águila. \\
        \midrule
        \multicolumn{5}{l}{\footnotesize \textit{Fuente:} Elaborado con traducciones propias a partir de \cite{diccionario_nahuatl2018}.} \\
        \bottomrule
    \end{tabularx}
\end{table}
Todos los pueblos comparten características etnográficas similares, pero orígenes completamente diferentes, siendo que
\termino{Oztotepec, Xicomulco} y \termino{Cuauhtenco} no son provenientes de la tribu \termino{Momoxca} de la región.
Hasta antes del siglo XIX, la zona comprendida como alcaldía de Milpa Alta no formaba parte de la capital, el primer intento
de volverlo parte de la demarcación ocurrió en el mandato del presidente José Antonio López de Santa Ana (sin fecha exacta),
pero ocurría un problema legislativo claro: La región estaba dividida por sus orígenes, en la región montañosa poniente - norte,
los pueblos originarios no se sentían parte de Milpa Alta ni de la alcaldía de Xochimilco, teniendo más en común con la alejada
alcaldía de Tlalpan, aunque estos a su vez proclamaban no tener nada que ver con ellos.

Este conflicto, es anterior al gobierno de Santa Ana y es conocido como \termino{El paraje Tulmiac} u \termino{Ojo de Agua}.
En el cual se descubrió un enorme ojo de agua natural proveniente de la zona
montañosa de \termino{La sierra del Chichinautzin (gran fumador en náhuatl)}, el cual dotaba de agua a la región
montañosa del sur. El problema es que esta no fue descubierta por los habitantes de \termino{Cuauhtenco, Xicomulco} o
\termino{Oztotepec} sino por pobladores del centro de \termino{Atocpan, Malacachtepec} y \termino{Tlacotenco}, haciendo que la
disputa escalara hasta llegar a los tribunales de la entonces intendencia del virreinato, donde el estado colonial determina
una prórroga para medir los parajes, pero esta solución quedó en el aire, ninguno de los dos bandos quedó conforme y el caso se
volvió a reabrir en el último gobierno de Santa Ana (1855). En él, el nuevo gobierno manda emisarios y con ello, nace el mito
de \termino{La virgen de la Asunción}, en el que se cuenta que un hombre al visitar el paraje del ojo de agua \termino{Tulmiac},
fue visitado por una virgen de aspecto blanco y celestial que le mencionó que el ojo de agua sería un regalo para la unión,
fraternidad y mutuo entendimiento de la región y que marcaría el inicio de la prosperidad para todos sus
\termino{descendientes o hijos}, dicha virgen llegaría en una representación material cuando se mandaría a hacer un
retrato en el municipio de Chalco de Díaz de Covarrubias, donde llegaría impregnada de \termino{Sabiduría y amor}. Con ello,
el gobierno federal incorpora a Milpa Alta como región federal dentro del entonces Distrito Federal. Conformando la alcaldía en
doce pueblos originarios (aunque en un inicio eran nueve). Las revueltas no fueron muchas y aunque la incertidumbre y la
incredulidad de los más férreos fue notoria en los primeros procesos, las dudas se disiparon cuando la \termino{Virgen de la Asunción}
llegó al centro de la recién formada Delegación de Milpa Alta. El mito prevalece hasta la fecha y los habitantes lo recuerdan
con extremo fervor. Pero la cohesión social nunca cuajó, pues hasta la fecha se tiene un mal estigma sobre los pueblos cercanos
a \termino{Xicomulco, Cuauhtenco, Oztotepec}; actual \termino{San Bartolomé Xicomulco, San Salvador Cuauhtenco y San Pablo Oztotepec}
respectivamente. Siendo que estos no son considerados en la toma de decisiones a nivel comunal; pues a los doce comuneros de la región
de Milpa Alta les basta con tener una mayoría de consensos, representada históricamente por los nueve pueblos originarios fundadores.

\textbf{¿De qué nos sirve este ejemplo?}, si lo dejamos como una historia meramente anecdótica, pierde todo su valor en este ensayo.
En \textbf{El laberinto de la Soledad}, Octavio Paz describe cómo la sociedad mexicana es una amalgama de culturas que no acaban de
convivir con el medio que les rodea y les separa (progresivamente) de un mutuo entendimiento. No es sorpresa de nadie que el gobierno
federal y el de la Ciudad de México haya cesado ya hace unos años por el entendimiento del mal llamado \termino{Gobierno de los Pueblos},
al comportarse con una resistencia que rebasa lo cívico y legalmente válido.

Por un lado, tenemos a los pueblos aislados dentro de una región como lo fue el caso anterior y por otro tenemos una realidad
más \termino{Citadina}. \termino{Gilles Deleuze}, con apoyo en las obras de \termino{Foucault}, define que estos \termino{Cismas}
o sucesos de cambio pueden denominarse \termino{dispositivos}, los cuales pueden ser culturales, ideológicos, políticos, económicos y,
más recientemente tecnológicos. A pesar de esto, notamos un cierto \termino{Estándar} en toda la cultura mexicana alejada u
\termino{originaria}; festividades fuertemente arraigadas a símbolos de la fe católica que combinan elementos de la cultura prehispánica,
la cultura colonial-española y el desarrollo post-revolucionario de sus regiones que son fuertemente defendidos no en una
legitimidad de los tres puntos anteriores, si no en una base imaginaria de \termino{Cultura milenaria}, sin saber exactamente de donde
proviene o que significa. Esta base es o fue defendida durante muchos años por la \termino{Generación Baby Boomer, X} y una parte creciente
de la \termino{generación Milenial y Generación Z}, pero poco a poco va cayendo en una despersonalización sistémica. A la mayor parte
de las generaciones actuales no les interesa saber el significado cultural de su región, su origen o sus figuras católica-religiosas,
les interesan aspectos más banales y de moda; no es raro encontrar una semejanza abismal en todas las \termino{Ferias de Pueblo} por
toda la república Mexicana: Venta de Comida rápida como alitas de pollo, \termino{Boneless}, \termino{Banderillas}, papas fritas o
\termino{Wafles}. Juegos de Feria con una seguridad precaria y altamente criticable y dos factores que nunca pueden faltar: La venta
sistemática de bebidas alcohólicas preparadas o de \termino{Descorche} durante todo lo que dure la fiesta y una incesante banda de música
\termino{Regional Mexicano}, interpretando temas de moda que se han popularizado en los últimos veinte años a causa de una cultura basada
en la apariencia económica y el estatus socio-cultural mayormente influenciado por los grupos delictivos conocidos como
\termino{Cárteles Mexicanos}. En este sentido, el fenómeno comercial del \termino{Regional Mexicano} funciona hoy como un término
paraguas que agrupa expresiones musicales profundamente heterogéneas, haciendo indispensable para el análisis sociológico separar
las composiciones de origen tradicional de aquellas que fungen como aparatos ideológicos de la hiper-mercantilización de la violencia:

\begin{itemize}
    \item \textbf{Temas de arraigo tradicional y religioso:} Históricamente, canciones como \textit{El Rey} o \textit{Caminos de Michoacán}
            fungían como crónicas de la vida rural, el honor y la cohesión social, mientras que cantos como \textit{La Guadalupana}
            reflejaban el profundo sincretismo católico del centro y sur del país. Estas composiciones fueron el pilar de consumo de
            las generaciones adultas, promoviendo una identidad comunitaria y una ``mexicanidad'' clásica.

    \item \textbf{Temas de exaltación ilícita y narcocultura:} En brutal contraste, subgéneros contemporáneos como los corridos
            bélicos y tumbados han desplazado la narrativa rural. Ya no operan como simples crónicas regionales, sino que exaltan
            un consumismo neoliberal extremo, normalizando las estructuras de poder paramilitar. Las letras que ostentan ropa de
            diseñador, vehículos de lujo y el tráfico de estupefacientes se han convertido en un mecanismo de validación social
            rápido ante la falta de movilidad económica formal.
\end{itemize}

Impulsada por la precarización laboral y las secuelas estructurales de la fragmentación social de las últimas dos décadas,
la narcocultura dejó de ser un fenómeno marginal para convertirse en el núcleo de la cultura pop nacional que ameniza los
espacios públicos. Como advierte el análisis de \cite{bravo2025narcocorridos}, ante un Estado ausente y la nula esperanza
de movilidad social por vías legítimas (como la educación o el empleo), el crimen organizado se erige para las juventudes
excluidas como un modelo aspiracional. Así, la festividad patronal pierde su propósito originario de comunión comunitaria
para convertirse en un escenario donde se glorifica y mercantiliza el estilo de vida criminal.

Este fenómeno cultural escapa a lo que Durkheim menciona en \textbf{Las Reglas del método sociológico}, donde la amalgamación de
ciertas conductas pueden entenderse como un todo y su influencia en la sociedad encamina hacia un fin común. En este caso, donde la
cultura de la apariencia, el consumo y la \termino{liquidez} o inmediatez del momento, revela que no se tiene un consenso
esperanzador de las nuevas generaciones (parte de la \termino{Generación Milenial} y \termino{Generación Z}) hacia su futuro,
es más común ahora encontrar personas de este grupo social, dividido por edad, que imitan a su \termino{influencer} de turno,
artista del momento o comentario que se \termino{viraliza} en las redes sociales. Lo que dificulta el nacimiento de un
pensamiento propio específico y personal.

En \cite{foucault2002vigilar}, Foucault explora la idea de un \termino{Dispositivo Panóptico}, aquel que permite el control de las masas
a través de un dominio absoluto del pensamiento, es decir la obediencia ciega y sus consecuencias inmediatas. Aunque el autor se
basa en un modelo de cárcel construida por el filósofo Jeremy Bentham en 1791 e implementada en otros países, donde la población o
\termino{Carcelarios} son constantemente vigiladas al punto de que si cometen una falta al orden establecido, serán castigados
con extrema inmediatez. Al paso del tiempo, estos aprenderán a que, aún sin saber que no están siendo vigilados, la consecuencia será
inmediata a sus acciones y otorgarán un control absoluto sobre su pensamiento y acción. La sociedad mexicana se comporta más parecido
a este espectro, teniendo claras diferencias entre estratos ideológicos y económicos pero con una meta clara: sobrevivir y recriminar
a quién no obedece las reglas que, aun sin saberlas; existen.
